% Section: Introduction
\section{Introduction}
The partition function $Z$ serves as the fundamental normalization constant for Boltzmann probability distributions, defining the thermal equilibrium states of atomic systems \cite{Kou2006}. However, evaluating $Z$ for high-dimensional systems is computationally intractable due to the vast Potential Energy Surface (PES) inherent in many-body problems \cite{Wang2024}. While the naive enumeration of microstates is limited to systems with fewer than 20 variables, realistic atomic configurations involve high-dimensional integrals that cannot be resolved through standard summation \cite{Christiansen2023}. In the Fe/MgO(001) system, this complexity is further compounded by a significant surface free energy mismatch—where the energy of bcc Fe ($\approx 2.9$ J/m²) substantially exceeds that of the MgO(001) substrate ($\approx 1.1$ J/m²)—driving the island (Volmer-Weber) mode despite a modest 3.8\% lattice mismatch \cite{Wisnios2015, Yuasa2004}. To address these limitations, an accurate characterization of $Z$ is essential, requiring a statistical ensemble that captures representative atomic morphologies.

Existing computational methods often fail to adequately sample these rugged PES. Traditional stochastic methods, such as Monte Carlo (MC) or Nested Sampling, require a prohibitive number of evaluations to achieve convergence when paired with first-principles Density Functional Theory (DFT) \cite{Hamamoto2023}. These samplers frequently become trapped in local minima, failing to resolve structural transitions critical for determining phase stability \cite{Noe2019}. Furthermore, generative machine learning architectures, including Normalizing Flows (NFs)\cite{papamakarios2021}, require extensive pre-computed datasets of $10^{3}$ to $10^{5}$ configurations, while variational methods remain inherently biased by the choice of trial distribution, often underestimating the true entropy of the system \cite{Paleico2020, Mehta2018}.

A limitation exists in the Boltzmann distribution, $P(x) \propto e^{-E(x)/k_{B} T}$, where $E(x)$ is the instantaneous energy of a configuration $x$, $k_{B}$ is the Boltzmann constant, and $T$ is the absolute temperature. This distribution cannot differentiate between local minima and transient states, such as saddle points or trajectory states \cite{david2017, angelani2002}. At a given energy level, local minima must be prioritized because they represent states of mechanical equilibrium, whereas trajectory states lack the finite lifetime required for physical manifestation \cite{somani2013, asenjo2013}. Under experimental conditions, atoms reside within basins of attraction where net forces are zero; thus, only these local minima represent the true polymorphs observable in experiment \cite{wales2006}. By restricting the Boltzmann factor to a "minimum-ensemble" of local minima, we can accurately represent the rugged PES of Fe/MgO(001), where island growth populates a landscape of competitive island structures.

In this work, by implementing the active learning framework with a biased exploration strategy, we successfully demonstrated partition function sampling, which represents a minimum ensemble of local minima, enabling the categorization of the PES into distinct island and flat growth modes of Fe on the MgO(001) substrate and revealing the temperature-dependent probability of each growth motif. This paper is organized as follows. In Sec. II, we present methods of the active learning framework to sample an ensemble of surrogate-relaxed structures. In Sec. III, we first discuss the GPR optimization and ground-state identification of Fe on MgO(001), and subsequently introduce the minimum ensemble of local minima. We summarize our work in Sec. IV.